% Options for packages loaded elsewhere
\PassOptionsToPackage{unicode}{hyperref}
\PassOptionsToPackage{hyphens}{url}
\documentclass[
]{article}
\usepackage{xcolor}
\usepackage[margin=1in]{geometry}
\usepackage{amsmath,amssymb}
\setcounter{secnumdepth}{-\maxdimen} % remove section numbering
\usepackage{iftex}
\ifPDFTeX
  \usepackage[T1]{fontenc}
  \usepackage[utf8]{inputenc}
  \usepackage{textcomp} % provide euro and other symbols
\else % if luatex or xetex
  \usepackage{unicode-math} % this also loads fontspec
  \defaultfontfeatures{Scale=MatchLowercase}
  \defaultfontfeatures[\rmfamily]{Ligatures=TeX,Scale=1}
\fi
\usepackage{lmodern}
\ifPDFTeX\else
  % xetex/luatex font selection
\fi
% Use upquote if available, for straight quotes in verbatim environments
\IfFileExists{upquote.sty}{\usepackage{upquote}}{}
\IfFileExists{microtype.sty}{% use microtype if available
  \usepackage[]{microtype}
  \UseMicrotypeSet[protrusion]{basicmath} % disable protrusion for tt fonts
}{}
\makeatletter
\@ifundefined{KOMAClassName}{% if non-KOMA class
  \IfFileExists{parskip.sty}{%
    \usepackage{parskip}
  }{% else
    \setlength{\parindent}{0pt}
    \setlength{\parskip}{6pt plus 2pt minus 1pt}}
}{% if KOMA class
  \KOMAoptions{parskip=half}}
\makeatother
\usepackage{color}
\usepackage{fancyvrb}
\newcommand{\VerbBar}{|}
\newcommand{\VERB}{\Verb[commandchars=\\\{\}]}
\DefineVerbatimEnvironment{Highlighting}{Verbatim}{commandchars=\\\{\}}
% Add ',fontsize=\small' for more characters per line
\usepackage{framed}
\definecolor{shadecolor}{RGB}{248,248,248}
\newenvironment{Shaded}{\begin{snugshade}}{\end{snugshade}}
\newcommand{\AlertTok}[1]{\textcolor[rgb]{0.94,0.16,0.16}{#1}}
\newcommand{\AnnotationTok}[1]{\textcolor[rgb]{0.56,0.35,0.01}{\textbf{\textit{#1}}}}
\newcommand{\AttributeTok}[1]{\textcolor[rgb]{0.13,0.29,0.53}{#1}}
\newcommand{\BaseNTok}[1]{\textcolor[rgb]{0.00,0.00,0.81}{#1}}
\newcommand{\BuiltInTok}[1]{#1}
\newcommand{\CharTok}[1]{\textcolor[rgb]{0.31,0.60,0.02}{#1}}
\newcommand{\CommentTok}[1]{\textcolor[rgb]{0.56,0.35,0.01}{\textit{#1}}}
\newcommand{\CommentVarTok}[1]{\textcolor[rgb]{0.56,0.35,0.01}{\textbf{\textit{#1}}}}
\newcommand{\ConstantTok}[1]{\textcolor[rgb]{0.56,0.35,0.01}{#1}}
\newcommand{\ControlFlowTok}[1]{\textcolor[rgb]{0.13,0.29,0.53}{\textbf{#1}}}
\newcommand{\DataTypeTok}[1]{\textcolor[rgb]{0.13,0.29,0.53}{#1}}
\newcommand{\DecValTok}[1]{\textcolor[rgb]{0.00,0.00,0.81}{#1}}
\newcommand{\DocumentationTok}[1]{\textcolor[rgb]{0.56,0.35,0.01}{\textbf{\textit{#1}}}}
\newcommand{\ErrorTok}[1]{\textcolor[rgb]{0.64,0.00,0.00}{\textbf{#1}}}
\newcommand{\ExtensionTok}[1]{#1}
\newcommand{\FloatTok}[1]{\textcolor[rgb]{0.00,0.00,0.81}{#1}}
\newcommand{\FunctionTok}[1]{\textcolor[rgb]{0.13,0.29,0.53}{\textbf{#1}}}
\newcommand{\ImportTok}[1]{#1}
\newcommand{\InformationTok}[1]{\textcolor[rgb]{0.56,0.35,0.01}{\textbf{\textit{#1}}}}
\newcommand{\KeywordTok}[1]{\textcolor[rgb]{0.13,0.29,0.53}{\textbf{#1}}}
\newcommand{\NormalTok}[1]{#1}
\newcommand{\OperatorTok}[1]{\textcolor[rgb]{0.81,0.36,0.00}{\textbf{#1}}}
\newcommand{\OtherTok}[1]{\textcolor[rgb]{0.56,0.35,0.01}{#1}}
\newcommand{\PreprocessorTok}[1]{\textcolor[rgb]{0.56,0.35,0.01}{\textit{#1}}}
\newcommand{\RegionMarkerTok}[1]{#1}
\newcommand{\SpecialCharTok}[1]{\textcolor[rgb]{0.81,0.36,0.00}{\textbf{#1}}}
\newcommand{\SpecialStringTok}[1]{\textcolor[rgb]{0.31,0.60,0.02}{#1}}
\newcommand{\StringTok}[1]{\textcolor[rgb]{0.31,0.60,0.02}{#1}}
\newcommand{\VariableTok}[1]{\textcolor[rgb]{0.00,0.00,0.00}{#1}}
\newcommand{\VerbatimStringTok}[1]{\textcolor[rgb]{0.31,0.60,0.02}{#1}}
\newcommand{\WarningTok}[1]{\textcolor[rgb]{0.56,0.35,0.01}{\textbf{\textit{#1}}}}
\usepackage{graphicx}
\makeatletter
\newsavebox\pandoc@box
\newcommand*\pandocbounded[1]{% scales image to fit in text height/width
  \sbox\pandoc@box{#1}%
  \Gscale@div\@tempa{\textheight}{\dimexpr\ht\pandoc@box+\dp\pandoc@box\relax}%
  \Gscale@div\@tempb{\linewidth}{\wd\pandoc@box}%
  \ifdim\@tempb\p@<\@tempa\p@\let\@tempa\@tempb\fi% select the smaller of both
  \ifdim\@tempa\p@<\p@\scalebox{\@tempa}{\usebox\pandoc@box}%
  \else\usebox{\pandoc@box}%
  \fi%
}
% Set default figure placement to htbp
\def\fps@figure{htbp}
\makeatother
\setlength{\emergencystretch}{3em} % prevent overfull lines
\providecommand{\tightlist}{%
  \setlength{\itemsep}{0pt}\setlength{\parskip}{0pt}}
\usepackage{bookmark}
\IfFileExists{xurl.sty}{\usepackage{xurl}}{} % add URL line breaks if available
\urlstyle{same}
\hypersetup{
  pdftitle={Bayesian Meta-Analysis},
  hidelinks,
  pdfcreator={LaTeX via pandoc}}

\title{Bayesian Meta-Analysis}
\author{}
\date{\vspace{-2.5em}2026-04-17}

\begin{document}
\maketitle

\subsection{Introduction}\label{introduction}

Meta-analysis with \(k\) studies is awkward for frequentist
random-effects estimators when \(k\) is small: REML and
DerSimonian-Laird both rely on asymptotic approximations that break down
when only five or six studies contribute, and the between-study variance
\(\tau^2\) can collapse to zero or be wildly imprecise. A Bayesian
formulation puts an explicit prior on \(\tau\) and propagates
uncertainty into both the pooled effect and the prediction interval for
a new study, which is exactly the inferential question that drives most
meta-analyses in clinical and biomedical work.

\subsection{Prerequisites}\label{prerequisites}

A working understanding of fixed- and random-effects meta-analysis,
exchangeability assumptions, and hierarchical Bayesian models.
Familiarity with \texttt{brms} formula syntax accelerates the
implementation step.

\subsection{Theory}\label{theory}

The hierarchical model is

\[y_i \sim \mathcal N(\theta_i, s_i^2), \qquad \theta_i \sim \mathcal N(\mu, \tau^2),\]

where \(y_i\) and \(s_i\) are the observed effect estimate and its
standard error from study \(i\), \(\theta_i\) is the latent
study-specific effect, \(\mu\) is the population effect, and \(\tau\) is
the between-study standard deviation. Typical priors are weakly
informative: \(\mu \sim \mathcal N(0, 2)\) on the
standardised-mean-difference scale and
\(\tau \sim \mathrm{half\text{-}Normal}(0, 0.5)\) or
\(\mathrm{half\text{-}Cauchy}(0, 0.3)\). The half-Cauchy gives the
heavier tail favoured when very large heterogeneity is plausible. The
posterior delivers \(\mu\), \(\tau\), every \(\theta_i\) shrunk toward
\(\mu\), and a predictive distribution \(\mathcal N(\mu, \tau^2)\) for
the next study.

\subsection{Assumptions}\label{assumptions}

Exchangeable studies, accurate within-study standard errors \(s_i\), and
proper priors. The within-study sampling distribution is taken as known;
if \(s_i\) are themselves estimated from few patients, an additional
level of hierarchy is needed.

\subsection{R Implementation}\label{r-implementation}

\begin{Shaded}
\begin{Highlighting}[]
\FunctionTok{library}\NormalTok{(brms)}

\NormalTok{studies }\OtherTok{\textless{}{-}} \FunctionTok{data.frame}\NormalTok{(}
  \AttributeTok{yi =} \FunctionTok{c}\NormalTok{(}\FloatTok{0.2}\NormalTok{, }\FloatTok{0.3}\NormalTok{, }\SpecialCharTok{{-}}\FloatTok{0.05}\NormalTok{, }\FloatTok{0.4}\NormalTok{, }\FloatTok{0.25}\NormalTok{),}
  \AttributeTok{sei =} \FunctionTok{c}\NormalTok{(}\FloatTok{0.10}\NormalTok{, }\FloatTok{0.15}\NormalTok{, }\FloatTok{0.08}\NormalTok{, }\FloatTok{0.12}\NormalTok{, }\FloatTok{0.20}\NormalTok{)}
\NormalTok{)}

\NormalTok{fit }\OtherTok{\textless{}{-}} \FunctionTok{brm}\NormalTok{(yi }\SpecialCharTok{|} \FunctionTok{se}\NormalTok{(sei) }\SpecialCharTok{\textasciitilde{}} \DecValTok{1} \SpecialCharTok{+}\NormalTok{ (}\DecValTok{1} \SpecialCharTok{|} \FunctionTok{seq\_len}\NormalTok{(}\FunctionTok{nrow}\NormalTok{(studies))),}
           \AttributeTok{data =}\NormalTok{ studies,}
           \AttributeTok{prior =} \FunctionTok{prior}\NormalTok{(}\FunctionTok{normal}\NormalTok{(}\DecValTok{0}\NormalTok{, }\DecValTok{1}\NormalTok{), }\AttributeTok{class =} \StringTok{"Intercept"}\NormalTok{) }\SpecialCharTok{+}
                   \FunctionTok{prior}\NormalTok{(}\FunctionTok{cauchy}\NormalTok{(}\DecValTok{0}\NormalTok{, }\FloatTok{0.5}\NormalTok{), }\AttributeTok{class =} \StringTok{"sd"}\NormalTok{),}
           \AttributeTok{chains =} \DecValTok{4}\NormalTok{, }\AttributeTok{iter =} \DecValTok{4000}\NormalTok{, }\AttributeTok{refresh =} \DecValTok{0}\NormalTok{)}

\FunctionTok{summary}\NormalTok{(fit)}
\end{Highlighting}
\end{Shaded}

\subsection{Output \& Results}\label{output-results}

\texttt{summary(fit)} returns the pooled effect (\texttt{Intercept}) and
the between-study SD (\texttt{sd(seq\_len(...))}) with 95 \% credible
intervals, \(\hat R\), and ESS. Per-study posterior means come from
\texttt{ranef(fit)}; the predictive distribution for a future study is
\texttt{posterior\_predict()} with \texttt{re\_formula\ =\ NA}.

\subsection{Interpretation}\label{interpretation}

A reporting sentence: ``A Bayesian random-effects meta-analysis with a
half-Cauchy(0, 0.5) prior on \(\tau\) pooled the five studies to an
effect of 0.23 (95 \% credible interval 0.07 to 0.39); the between-study
heterogeneity was small (\(\tau = 0.09\), 95 \% CrI 0.01 to 0.25), and
the 95 \% predictive interval for a future trial was \(-0.05\) to
\(0.51\).'' Reporting the predictive interval together with the pooled
estimate avoids the false sense of precision that a credible interval on
\(\mu\) alone can convey.

\subsection{Practical Tips}\label{practical-tips}

\begin{itemize}
\tightlist
\item
  With \(k < 10\) studies, the prior on \(\tau\) matters substantially;
  a half-\(t(3, 0, 1)\) or half-Normal(0, 0.5) is a reasonable
  compromise between letting the data speak and avoiding implausible
  values.
\item
  Always report the prior in the methods section and conduct a
  sensitivity analysis with at least one wider and one tighter prior on
  \(\tau\).
\item
  For binary outcomes, fit on the log-odds-ratio or risk-ratio scale
  rather than risk differences; the Normal-Normal hierarchy is more
  reasonable on the log scale.
\item
  The \texttt{bayesmeta} package provides closed-form computations for
  simple models and is faster than MCMC when the prior is conjugate.
\item
  For network meta-analysis with multiple treatments, switch to
  \texttt{gemtc} or \texttt{multinma}; the random-effects machinery
  generalises but the implementation is non-trivial.
\end{itemize}

\subsection{R Packages Used}\label{r-packages-used}

\texttt{brms} for the hierarchical Bayesian fit with study-level
standard errors, \texttt{bayesmeta} for closed-form alternatives,
\texttt{bayesplot} for diagnostic plots, and \texttt{gemtc} /
\texttt{multinma} for the network-meta extension.

\subsection{For Reviewers}\label{for-reviewers}

\textbf{What to look for in a paper using this method.}

\begin{itemize}
\tightlist
\item
  \textbf{Common misapplications.}
\item
  \textbf{Diagnostics that should be reported but often aren't.}
\item
  \textbf{Red flags in tables and figures.}
\item
  \textbf{What to verify.}
\item
  \textbf{What an adequate Methods paragraph must contain.}
\end{itemize}

\subsection{See also --- labs in this
chapter}\label{see-also-labs-in-this-chapter}

\begin{itemize}
\tightlist
\item
  \href{labs/lab-bayesian-thinking-control-vs-decision-error.Rmd}{Bayesian
  thinking; α control vs decision error}
\item
  \href{labs/lab-brms-stan-loo-hierarchical-models.Rmd}{brms / Stan,
  LOO, hierarchical models}
\end{itemize}

\end{document}
